\section{A translation calculus for Hurwitz jets}
\label{sec:translation}

The parameter identity $\partial_x\zeta(s,x)=-s\zeta(s+1,x)$
connects the three central objects of this section: Hurwitz order derivatives,
Stieltjes antiderivatives, and the logarithm of the Gamma function.
We add translation in $x$ as a second operation. The resulting bilinear
identity has both a polylogarithmic form and a Hurwitz form. The latter
makes all endpoint logarithms and all derivative orders explicit.

\subsection{Normalization and analytic justification}

Write
\begin{equation}\label{eq:stieltjes-def}
 \zeta(1+u,x)=\frac1u+\sum_{n=0}^{\infty}
              \frac{(-1)^n}{n!}\gamma_n(x)u^n,
 \qquad x>0.
\end{equation}
Order derivatives of $\zeta$ always refer to its first argument. Define
\begin{equation}\label{eq:FJ}
 F_r(x)=\zeta^{(r)}(0,x),\qquad
 J_n(x)=\frac{(-1)^{n+1}}{n+1}F_{n+1}(x)
 \quad(r,n\geq0).
\end{equation}
Then
\[
 F_0(x)=\frac12-x,\qquad
 F_1(x)=\log\Gamma(x)-\frac12\log(2\pi),\qquad J_n'(x)=\gamma_n(x).
\]
The last equality follows by comparing coefficients in
$\partial_x\zeta(s,x)=-s\zeta(s+1,x)$ at $s=0$. This classical ladder is
already present in the manuscript; its earlier literature is discussed
there. Our normalization is the mean-zero one:
\begin{equation}\label{eq:mean-zero}
 \int_0^1F_r(x)\,dx=\int_0^1J_n(x)\,dx=0.
\end{equation}
Thus $J_n$ differs by a specified constant from the manuscript's primitive
normalized to vanish at $1$.

All functions in \eqref{eq:FJ} are regarded as periodic functions almost
everywhere on $\mathbb T=\mathbb R/\mathbb Z$. Endpoint values themselves
are immaterial. Put
\[
 \Delta_hf(x)=f(\{x+h\})-f(x),\qquad 0<h<1.
\]
This is a periodic translation, not a translation of the original
nonperiodic Gamma function on the whole positive half-line.

For completeness, the analytic justification for all forthcoming
coefficient operations can be made at the level of Banach-space-valued
holomorphic functions. On every disk $|s|\leq\rho<1/2$,
\[
 \zeta(s,x)=x^{-s}+\zeta(s,1+x),\qquad 0<x<1,
\]
and the second term and its order derivatives are uniformly bounded on
smaller disks. Consequently the map $s\mapsto\zeta(s,\cdot)$ is
holomorphic with values in $L^2(0,1)$; its $r$th derivative is dominated
by a constant times $1+x^{-\rho}|\log x|^r$. Translation is a bounded
operator on $L^2$. The bilinear products below are therefore jointly
holomorphic, and order differentiation under the integrals is valid.
The identity $\int_0^1\zeta(s,x)\,dx=0$ for $\Re s<1$ proves
\eqref{eq:mean-zero} by the same argument.

\subsection{The exact master identity}

For $u=s+t$ set
\begin{align}
 \mathcal R(s,t)&=
 \frac{\Gamma(1-s)\Gamma(1-t)}{\Gamma(2-s-t)}
 \frac{\cos\bigl(\pi(s-t)/2\bigr)}{\cos\bigl(\pi(s+t)/2\bigr)},
 \label{eq:R-kernel}\\
 \mathcal Z(u,h)&=\zeta(u-1,h)+\zeta(u-1,1-h)-2\zeta(u-1).
 \label{eq:Z-kernel}
\end{align}
Both expressions are holomorphic near $(s,t)=(0,0)$.

\begin{theorem}[Exact translation identity]\label{thm:translation}
For $0<h<1$ and $|s|,|t|<1/4$,
\begin{equation}\label{eq:translation-master}
 \boxed{\quad
 \int_0^1\Delta_h\zeta(s,x)\,\Delta_h\zeta(t,x)\,dx
      =\mathcal R(s,t)\mathcal Z(s+t,h).
 \quad}
\end{equation}
In particular, for every $r,q\geq0$,
\begin{equation}\label{eq:translation-jets}
 \int_0^1\Delta_hF_r(x)\Delta_hF_q(x)\,dx
 =\left.\partial_s^r\partial_t^q
        \bigl(\mathcal R(s,t)\mathcal Z(s+t,h)\bigr)\right|_{s=t=0}.
\end{equation}
Multiplication by $(-1)^{m+n+2}/((m+1)(n+1))$ gives the corresponding
identity for $J_m,J_n$ when $r=m+1$, $q=n+1$.
\end{theorem}

\begin{proof}
Use the Fourier convention $\widehat f(k)=\int_0^1 f(x)e^{-2\pi ikx}\,dx$.
Hurwitz's Fourier formula gives
\begin{equation}\label{eq:fourier-hurwitz}
 \widehat{\zeta(s,\cdot)}(0)=0,
 \qquad
 \widehat{\zeta(s,\cdot)}(k)=\Gamma(1-s)(2\pi ik)^{s-1}
 \quad(k\ne0),
\end{equation}
where $\log(2\pi ik)=\log(2\pi|k|)+\tfrac{\pi i}2\sgn k$.
Initially one may take $\Re s<0$ and then extend in $L^2$.
These are classical Fourier coefficients \cite{DLMFHurwitz,EM1}.
Parseval's identity for the bilinear pairing, together with
$\Delta_h$ acting by $e^{2\pi ikh}-1$, yields
\begin{align}
 \mathcal D(s,t;h)
 &=4\Gamma(1-s)\Gamma(1-t)(2\pi)^{u-2}
       \cos\frac{\pi(s-t)}2\notag\\
 &\quad\cdot\left\{\zeta(2-u)
       -\frac{\Li_{2-u}(e^{2\pi ih})+
                    \Li_{2-u}(e^{-2\pi ih})}{2}\right\}.
 \label{eq:polylog-master}
\end{align}
The series on this line are absolutely and locally uniformly convergent
in the stated bidisk. Notice that the symmetric pair of polylogarithms
must be retained when $s,t$ are complex; writing a real part would then
destroy holomorphic dependence on the orders.

The even part of the same Hurwitz Fourier formula, now at order $u-1$,
is
\begin{align*}
 \zeta(u-1,h)+\zeta(u-1,1-h)
 &=-4\Gamma(2-u)(2\pi)^{u-2}\cos\frac{\pi u}{2}\,
      \frac{\Li_{2-u}(e^{2\pi ih})+\Li_{2-u}(e^{-2\pi ih})}{2},\\
 2\zeta(u-1)&=-4\Gamma(2-u)(2\pi)^{u-2}
                   \cos\frac{\pi u}{2}\,\zeta(2-u).
\end{align*}
Their difference converts \eqref{eq:polylog-master} into
\eqref{eq:translation-master}. The $L^2$ argument above proves the
differentiated identity.
\end{proof}

\begin{remark}[Relation to the classical product integral]
At zero shift the underlying correlation is the classical identity
\[
 \int_0^1\zeta(s,x)\zeta(t,x)\,dx
 =-\mathcal R(s,t)\zeta(s+t-1),
\]
given, for example, by Espinosa--Moll \cite[Theorem 3.1]{EM1}.
We do not claim this antecedent as new. The contribution used here is
the translated difference, its uniformly justified jet calculus, and the
explicit consequences below. No claim of global priority is made for
every special-case identity obtained by differentiating the master formula.
\end{remark}

\subsection{Convergent endpoint expansion at every order}

Put $\ell=\log(1/h)$. The recurrence for Hurwitz zeta gives
\begin{equation}\label{eq:Z-split}
 \mathcal Z(u,h)=h^{1-u}+\zeta(u-1,1+h)+\zeta(u-1,1-h)-2\zeta(u-1).
\end{equation}
The last three terms are analytic and even in $h$ for $|h|<1$.
Expanding them around $h=0$ gives the particularly useful formula
\begin{align}
 \mathcal D(s,t;h)=\mathcal R(s,t)\biggl\{
 &h^{1-u}+u(u-1)\zeta(1+u)h^2\notag\\
 &+2\sum_{j=2}^{\infty}
   \frac{(u-1)_{2j}}{(2j)!}\zeta(2j-1+u)h^{2j}\biggr\}.
 \label{eq:convergent-expansion}
\end{align}
The product $u(u-1)\zeta(1+u)$ has its removable value $-1$ at zero.
This cancellation is performed before differentiation.

\begin{theorem}[All-order logarithmic endpoint polynomial]
\label{thm:cusp}
For fixed $r,q\geq0$ there is a monic polynomial of degree $r+q$,
\begin{equation}\label{eq:P-def}
 P_{r,q}(\ell)=
 \left.\partial_s^r\partial_t^q
     \bigl(\mathcal R(s,t)e^{(s+t)\ell}\bigr)\right|_{s=t=0},
\end{equation}
such that
\begin{equation}\label{eq:cusp}
 \int_0^1\Delta_hF_r\,\Delta_hF_q\,dx
   =hP_{r,q}(\log(1/h))+
           \sum_{j=1}^{\infty}a_{r,q,j}h^{2j},\qquad 0<h<1.
\end{equation}
The even series is convergent, locally uniformly for $|h|<1$.
If $r+q\geq1$, the coefficient of $\ell^{r+q-1}$ is $r+q$.
Every coefficient of $P_{r,q}$ lies in the algebra over $\mathbb Q$
generated by $\zeta(2),\ldots,\zeta(r+q)$; Euler's constant cancels.
The coefficient $a_{r,q,1}$ uses only Stieltjes constants through
$\gamma_{r+q-1}$ and these zeta values. For $j\geq2$, it uses
derivatives of $\zeta$ at $2j-1$ of order at most $r+q-1$.
For $r=q=0$ the statements about negative derivative indices mean that
only the constant $a_{0,0,1}=-1$ is present.
\end{theorem}

\begin{proof}
Equation \eqref{eq:convergent-expansion}, differentiated in the bidisk,
gives \eqref{eq:cusp}. Uniform convergence with order derivatives follows
also directly from the analytic Taylor expansion in \eqref{eq:Z-split}
on any smaller circle $|h|\leq h_0<1$.
The required coefficients of $\mathcal R$ are especially transparent in
its logarithm:
\begin{align}
 \log\mathcal R(s,t)
 &=\sum_{k\geq1}\frac{(s+t)^k}{k}
   +\sum_{k\geq2}\frac{\zeta(k)}k
      \bigl(s^k+t^k-(s+t)^k\bigr)\notag\\
 &\quad+\sum_{m\geq1}\frac{(1-2^{-2m})\zeta(2m)}m
      \bigl((s+t)^{2m}-(s-t)^{2m}\bigr).
 \label{eq:logR}
\end{align}
This follows from $\Gamma(2-u)=(1-u)\Gamma(1-u)$ and the convergent
logarithmic Gamma and cosine products. In particular,
$\mathcal R(0,0)=1$ and $\partial_s\mathcal R(0,0)
=\partial_t\mathcal R(0,0)=1$. These facts give the two leading
coefficients of $P$. The same formula proves the claimed coefficient
algebra. Finally $(u-1)_{2j}$ has a simple zero at $u=0$ for every
$j\geq1$. At $j=1$ it removes the pole of $\zeta(1+u)$; at $j\geq2$
it lowers the highest needed zeta derivative by one.
\end{proof}

Here are the first nontrivial endpoint polynomials:
\begin{align}
 P_{1,1}(\ell)&=\ell^2+2\ell+2+\frac{\pi^2}{3},\label{eq:P11}\\
 P_{1,2}(\ell)&=\ell^3+3\ell^2+
       \left(6+\frac{2\pi^2}{3}\right)\ell
       +6+\frac{2\pi^2}{3}-2\zeta(3),\label{eq:P12}\\
 P_{2,2}(\ell)&=\ell^4+4\ell^3+
       \left(12+\frac{4\pi^2}{3}\right)\ell^2
       +\left(24+\frac{8\pi^2}{3}-8\zeta(3)\right)\ell\notag\\
 &\quad+24+\frac{8\pi^2}{3}-8\zeta(3)+\frac{7\pi^4}{45}.
 \label{eq:P22}
\end{align}
The supplied exact-arithmetic program constructs these polynomials and
additional ones from \eqref{eq:logR}. Formula \eqref{eq:cusp} is stronger
than a finite asymptotic expansion: after the single logarithmic term,
the remaining function is an ordinary convergent even power series.

\subsection{A positive integral for the endpoint kernel}

The Gamma--trigonometric factor has an elementary positive-kernel
interpretation, which gives information about all the endpoint
polynomials simultaneously.

\begin{theorem}[Endpoint Gram identity]\label{thm:endpoint-gram}
For $|s|,|t|<1/4$, with $u=s+t$,
\begin{equation}\label{eq:R-integral}
 \mathcal R(s,t)=\frac1{1-u}+
 \int_0^\infty
 [(x+1)^{-s}-x^{-s}][(x+1)^{-t}-x^{-t}]\,dx.
\end{equation}
For every real $\ell$ this implies
\begin{align}
 P_{r,q}(\ell)={}&\int_0^1(\ell-\log x)^{r+q}\,dx\notag\\
 &+\int_0^\infty
 \bigl[(\ell-\log(x+1))^r-(\ell-\log x)^r\bigr]\notag\\
 &\hspace{18mm}\cdot
 \bigl[(\ell-\log(x+1))^q-(\ell-\log x)^q\bigr]\,dx.
 \label{eq:P-gram}
\end{align}
Every finite principal matrix of $(P_{r,q}(\ell))_{r,q\geq0}$ is
positive definite. For the initial blocks
$P_N(\ell)=(P_{r,q}(\ell))_{0\leq r,q\leq N}$,
\begin{equation}\label{eq:P-determinant}
 \det P_N(\ell)=\det P_N(0)\geq\prod_{r=0}^N(r!)^2,
\end{equation}
with strict inequality for $N\geq1$.
\end{theorem}

\begin{proof}
The integral in \eqref{eq:R-integral} is locally holomorphic: at
infinity its integrand is $O(x^{-\Re u-2})$, with fixed logarithmic
factors after differentiation, and at zero the factors are bounded
by a constant times $1+x^{-\rho}$ for any fixed $\rho<1/2$ on a
smaller bidisk. Make the substitution $y=x/(1+x)$. The integral
becomes
\[
 \int_0^1(1-y)^{u-2}
       [1-y^{-s}-y^{-t}+y^{-u}]\,dy.
\]
The bracket has a double zero at $y=1$. To evaluate it without
separating divergent integrals, use the twice-subtracted beta formula
\begin{align*}
 B(a,b)={}&\int_0^1(1-y)^{b-1}
       [y^{a-1}-1+(a-1)(1-y)]\,dy\\
 &+\frac1b-\frac{a-1}{b+1}.
\end{align*}
It follows from the ordinary beta integral when $\Re b>0$ and
continues to $\Re a>0$, $\Re b>-2$, away from its displayed poles.
Apply it to $a=1,1-s,1-t,1-u$ with signs $+,-,-,+$ and
$b=u-1$. Both subtraction terms cancel. Thus the original integral
equals
\[
 -\frac1{1-u}-\Gamma(u-1)
 \left[\frac{\Gamma(1-s)}{\Gamma(t)}
             +\frac{\Gamma(1-t)}{\Gamma(s)}\right].
\]
The term $B(1-u,u-1)$ is zero under this meromorphic continuation.
Reflection of the Gamma factors, followed by
\[
 \frac{\sin\pi s+\sin\pi t}{\sin\pi(s+t)}
 =\frac{\cos(\pi(s-t)/2)}{\cos(\pi(s+t)/2)},
\]
identifies the last expression with $\mathcal R(s,t)-(1-u)^{-1}$.
All apparent singularities at $u=0$ are removable in the combined
identity, proving \eqref{eq:R-integral} throughout the bidisk.

Multiply by $e^{u\ell}$ and differentiate at zero to obtain
\eqref{eq:P-gram}. The two terms are Gram matrices of real functions.
The first is strictly positive definite for any distinct finite
set of indices: a nonzero polynomial in $\ell-\log x$ cannot
vanish almost everywhere on $(0,1)$.

Let $T(\ell)_{r,j}=\binom rj\ell^{r-j}$ for $j\leq r$ and zero
otherwise. The binomial theorem gives the congruence
$P_N(\ell)=T(\ell)P_N(0)T(\ell)^{\mathsf T}$ and
$\det T(\ell)=1$. For the first Gram term alone the matrix at
$\ell=0$ is $((r+q)!)$. Its determinant is
$\prod_{r=0}^N(r!)^2$, because
\[
 \binom{r+q}{r}=\sum_k\binom rk\binom qk
\]
is a factorization by a unit lower triangular Pascal matrix and its
transpose. Adding a positive semidefinite matrix cannot decrease
the determinant of a positive definite matrix. The second term of
\eqref{eq:P-gram} is nonzero for $N\geq1$, as its $(1,1)$ entry is
the integral of $\log^2(1+1/x)>0$. Hence the determinant increases
strictly. This proves \eqref{eq:P-determinant}. The shift-invariance
assertion concerns the initial blocks; arbitrary principal submatrices
need not have shift-invariant determinants.
\end{proof}

For example, $\det P_1(\ell)=1+\pi^2/3$. The $2\times2$ case of
the Gram positivity gives $P_{r,q}(\ell)^2<P_{r,r}(\ell)P_{q,q}(\ell)$
for distinct $r,q$. These inequalities hold for every real $\ell$,
not merely for the large positive logarithms arising in a small shift.

\subsection{Higher antiderivatives and sharp regularity}

For $m\geq0$ define
\begin{equation}\label{eq:higher-primitive}
 F_{r,m}(x)=\left.\partial_s^r
       \frac{\zeta(s-m,x)}{(1-s)_m}\right|_{s=0},
 \qquad (1-s)_0=1.
\end{equation}
The identity $\partial_x\zeta(s-m,x)=(m-s)\zeta(s-m+1,x)$
shows $\partial_xF_{r,m}=F_{r,m-1}$ for $m\geq1$.
These are the uniquely normalized successive periodic primitives: every
one has mean zero. In particular
\[
 \frac{(-1)^{n+1}}{n+1}F_{n+1,m-1}
 \quad(m\geq1)
\]
is an $m$-fold antiderivative of $\gamma_n$ on $(0,1)$.
The mean-zero convention removes all polynomial ambiguities.

\begin{theorem}[Exact Sobolev thresholds]\label{thm:sobolev}
For each $r,m\geq0$ and real $\sigma$,
\[
 F_{r,m}\in H^\sigma(\mathbb T)
 \quad\Longleftrightarrow\quad \sigma<m+\frac12.
\]
Thus the $m$-fold antiderivative of $\gamma_n$, with $m\geq1$, has
threshold $m-1/2$. Moreover
\begin{equation}\label{eq:J-sharp}
 \|\Delta_hJ_n\|_2^2
 \sim\frac{h\log^{2n+2}(1/h)}{(n+1)^2}\qquad(h\downarrow0).
\end{equation}
For any nonzero finite combination $G=\sum_{n=0}^Nc_nJ_n$ with
$c_N\ne0$, the right side becomes
$c_N^2h\log^{2N+2}(1/h)/(N+1)^2$ when the coefficients are real.
\end{theorem}

\begin{proof}
Let
\[
 Q_r(\lambda)=\left.\partial_s^r
                      \bigl(\Gamma(1-s)e^{\lambda s}\bigr)\right|_{s=0}.
\]
This is a monic polynomial in $\lambda$ of degree $r$; for example
$Q_1(\lambda)=\lambda+\gamma$ and
$Q_2(\lambda)=(\lambda+\gamma)^2+\zeta(2)$.
Using $\Gamma(m+1-s)/(1-s)_m=\Gamma(1-s)$ in
\eqref{eq:fourier-hurwitz},
\begin{equation}\label{eq:primitive-fourier}
 \widehat{F_{r,m}}(k)
       =\frac{Q_r(\log(2\pi ik))}{(2\pi ik)^{m+1}}
          \quad(k\ne0).
\end{equation}
The squared $H^\sigma$ norm is therefore governed, up to positive
constant factors for large $k$, by
$\sum_{k\geq2}k^{2\sigma-2m-2}\log^{2r}k$.
The integral test proves precisely the stated threshold, including
divergence at equality. Equation \eqref{eq:J-sharp} follows from the
monic polynomial in Theorem~\ref{thm:cusp}; in a finite combination its
highest logarithmic degree cannot cancel.
\end{proof}

Every finite Gram matrix
\[ (\int_0^1F_r(x)F_q(x)\,dx)_{0\leq r,q\leq N} \]
is also strictly positive definite. If a linear
combination vanished in $L^2$, its Fourier coefficients would make a
polynomial in $\log(2\pi k)+i\pi/2$ vanish for all positive integers
$k$. That polynomial is identically zero; triangularity of the monic
$Q_r$ then forces every coefficient of the combination to vanish.
This is functional linear independence in $L^2$, not an assertion of
arithmetic independence of special values.

\subsection{Antiderivatives with respect to the translation}

There is a second antiderivative operation: integrate the translated
pairing in its shift, rather than integrate each factor in its argument.
It also remains inside the Hurwitz hierarchy. For $m\geq1$ put
\begin{align}
 \mathcal A_m(s,t;h)=\mathcal R(s,t)\biggl\{
 &\frac{\zeta(u-1-m,h)+(-1)^m\zeta(u-1-m,1-h)}{(2-u)_m}
 \notag\\[-2pt]
 &-\frac{2\zeta(u-1)h^m}{m!}\biggr\},\qquad u=s+t.
 \label{eq:shift-primitive}
\end{align}

\begin{corollary}[The complete shift-antiderivative ladder]
\label{cor:shift-primitive}
For $|s|,|t|<1/4$ and $0<h<1$,
$\partial_h^m\mathcal A_m(s,t;h)=\mathcal D(s,t;h)$.
The repeated definite integral is
\begin{align}
 &\frac1{(m-1)!}\int_0^h(h-v)^{m-1}\mathcal D(s,t;v)\,dv
 =\mathcal A_m(s,t;h)\notag\\
 &\hspace{5mm}-2\mathcal R(s,t)
 \sum_{k=1}^{\lfloor m/2\rfloor}
 \frac{\zeta(u-1-2k)}{(2-u)_{2k}}
 \frac{h^{m-2k}}{(m-2k)!}.
 \label{eq:shift-definite}
\end{align}
Every finite derivative in $s,t$ may be taken in this identity.
\end{corollary}

\begin{proof}
The parameter recurrence gives
$\partial_h^m\zeta(u-1-m,h)=(2-u)_m\zeta(u-1,h)$.
The reflected term acquires the additional factor $(-1)^m$, which
cancels its displayed sign. Differentiating the polynomial term gives
$-2\zeta(u-1)$, proving the indefinite identity.
To normalize, use the Hurwitz shift formula as $h\downarrow0$.
The term $h^{m+1-u}$ and its first $m-1$ derivatives vanish, and
the remaining values at zero are
\[
 \partial_h^j\mathcal A_m(s,t;0)
 =\mathcal R(s,t)
 \frac{[1+(-1)^{m-j}]\zeta(u-1-m+j)}{(2-u)_{m-j}},
 \qquad 0\leq j<m.
\]
Subtracting their Taylor polynomial gives \eqref{eq:shift-definite}.
The same endpoint bound with fixed powers of $\log h$ justifies all
order derivatives. In particular the sum is empty for $m=1$.
\end{proof}
